\chapter{Anotaciones}

\begin{itemize}
    \item Como en las prácticas anteriores, cuando se utilizaba el mismo caso de uso en diferentes situaciones, lo distinguíamos mediante colores para indicar dónde se declaraba y dónde se utilizaba. En esta práctica seguimos el mismo enfoque con los casos de uso ``Reembolso Bizum'', ``Reembolso PayPal'' y ``Reembolso Tarjeta''. Solo desarrollamos el caso de uso en el paquete donde consideramos que se declara.
    \item Algunas tablas denominadas ``Otros Datos'' y ``Comentarios'' están vacías, ya que, como se ha explicado en clase, contienen información redundante.
\end{itemize}